%!TeX root=../thesis.tex

% Keywords must be defined before the abstract/resumo commands.
\keywords{Software Architecture, Architectural Uncertainty, Hypotheses Engineering, ArchHypo, Hypo-Stage, Internal Developer Portal, Backstage, Case Study, Empirical Evaluation}

\palavraschave{Arquitetura de Software, Incerteza Arquitetural, Engenharia de Hipóteses, ArchHypo, Hypo-Stage, Portal do Desenvolvedor Interno, Backstage, Estudo de Caso, Avaliação Empírica}

% English abstract (primary)
\abstract{
Software engineering teams in continuous-delivery organizations routinely face
architectural decisions under incomplete information, yet few lightweight tools
exist to help them record, prioritize, and learn from those uncertainties
collaboratively. The ArchHypo technique applies hypotheses engineering to
software architecture: teams formulate architectural uncertainties as
falsifiable hypotheses, assess each by uncertainty level and impact, and
define a technical plan for handling it. Prior evaluations demonstrated the
technique's value but reported a steep learning curve and dependence on expert
facilitation, both exacerbated by the absence of tool support.

This dissertation presents the first empirical account of a tool-mediated
application of ArchHypo in a professional software development organization. The artifact is Hypo-Stage,
an open-source Backstage plugin deployed on a large organization's Internal
Developer Portal, where engineers already work. The study adopts a flexible,
multiple case-study design with formative evaluation and action-research
elements, following the framework of \citet{robson2016}. Two teams --- a
Product team of twelve engineers and a Platform team of four --- used Hypo-Stage
over two agile sprints totalling four weeks. Data were collected through the
author's participant observation (continuous field notes), artifact analysis of
the thirteen hypotheses and associated technical plans registered in Hypo-Stage,
and open-ended active listening sessions with the senior technical leaders of
each team. Analysis followed the three-stage thematic coding procedure
described by \citet{robson2016} for qualitative data in flexible designs.

Four contributions emerge from the work.
\textbf{(1)} A detailed observational account of how two teams in a large-scale
product organization adopted and used Hypo-Stage, characterizing hypothesis
registration dynamics, role stratification, and tool adoption patterns.
\textbf{(2)} A characterization of the architectural uncertainties the teams
made explicit: their sources, quality attributes (dominated by Reliability and
Auditability), uncertainty-impact profiles, and technical plan strategies.
\textbf{(3)} Five novel technical architectural hypothesis archetypes ---
\textit{Hypothesis Cluster}, \textit{Regulatory Trigger Hypothesis},
\textit{Decoupling Seam Hypothesis}, \textit{Load-Validated Hypothesis}, and
\textit{Shared Resource Governance Hypothesis} --- derived from thematic analysis
of the hypothesis dataset. These are the first empirically grounded
characterization of recurring hypothesis shapes in the ArchHypo research line,
directly addressing a gap identified in prior work; they are presented as
emerging patterns requiring further validation.
\textbf{(4)} Nine categories of empirically grounded refinements to Hypo-Stage,
each traced to a concrete observation and linked to a code commit, including a
hypothesis evolution chart, a prioritization dashboard, and inline formulation
guidance that reduces the technique's learning barrier.
}

% Portuguese abstract (required by IME-USP regulations)
\resumo{
Equipes de engenharia de software em organizações de entrega contínua tomam
rotineiramente decisões arquiteturais sob informação incompleta, mas dispõem de
poucas ferramentas leves para registrar, priorizar e aprender com essas
incertezas de forma colaborativa. A técnica ArchHypo aplica a engenharia de
hipóteses à arquitetura de software: as equipes formulam incertezas
arquiteturais como hipóteses falsificáveis, avaliam cada uma por nível de
incerteza e impacto, e definem um plano técnico para tratá-las. Avaliações
anteriores demonstraram o valor da técnica, mas reportaram uma curva de
aprendizado acentuada e dependência de facilitação especializada, ambos
agravados pela ausência de suporte ferramental.

Esta dissertação apresenta o primeiro relato empírico de uma aplicação mediada
por ferramenta do ArchHypo em uma organização profissional de desenvolvimento de software. O artefato é o Hypo-Stage,
um plugin \textit{open source} para o Backstage, implantado no Portal do
Desenvolvedor Interno de uma grande organização de software, onde os engenheiros
já atuam cotidianamente. O estudo adota um desenho flexível de múltiplos
estudos de caso, com propósito avaliativo formativo e elementos de
pesquisa-ação, seguindo o arcabouço de \citet{robson2016}. Duas equipes --- uma
equipe de Produto com doze engenheiros e uma equipe de Plataforma com quatro ---
utilizaram o Hypo-Stage ao longo de duas \textit{sprints} ágeis, totalizando
quatro semanas. Os dados foram coletados por meio de observação participante
do próprio autor (notas de campo contínuas), análise de artefatos das treze
hipóteses e planos técnicos registrados no Hypo-Stage, e sessões abertas de
escuta ativa com os líderes técnicos seniores de cada equipe. A análise
seguiu o procedimento de codificação temática em três estágios descrito por
\citet{robson2016} para dados qualitativos em designs flexíveis.

Quatro contribuições emergem do trabalho.
\textbf{(1)} Um relato observacional detalhado de como duas equipes em uma
organização de produto em escala adotaram e utilizaram o Hypo-Stage,
caracterizando a dinâmica de registro de hipóteses, a estratificação por papel
e os padrões de adoção da ferramenta.
\textbf{(2)} Uma caracterização das incertezas arquiteturais tornadas explícitas
pelas equipes: suas fontes, atributos de qualidade (dominados por
Confiabilidade e Auditabilidade), perfis de incerteza-impacto e estratégias
de plano técnico.
\textbf{(3)} Cinco arquétipos técnicos de hipóteses arquiteturais ---
\textit{Hypothesis Cluster}, \textit{Regulatory Trigger Hypothesis},
\textit{Decoupling Seam Hypothesis}, \textit{Load-Validated Hypothesis} e
\textit{Shared Resource Governance Hypothesis} --- derivados da análise temática
do conjunto de hipóteses. Estes constituem a primeira caracterização empírica de
formas recorrentes de hipóteses na linha de pesquisa do ArchHypo, abordando
diretamente uma lacuna identificada em trabalhos anteriores; são apresentados
como padrões emergentes que requerem validação adicional.
\textbf{(4)} Nove categorias de refinamentos empiricamente fundamentados ao
Hypo-Stage, cada um rastreado a uma observação concreta e vinculado a um
\textit{commit} no repositório, incluindo um gráfico de evolução das hipóteses,
um painel de priorização e orientações de formulação \textit{inline} que
reduzem a barreira de aprendizado da técnica.
}
